\section{Core Code Owns Every Decision an Optimizer Could Game}
\label{sec:v2:evaluation}

An optimizer that can read the answers, the scorer or the final set can look better than it is, and one that scores itself will bank noise. \capevolve{} therefore moves every such decision into core code that the optimizer cannot edit. Splits are seeded, written once and read back by every skill. The final split is scored once: the run directory raises an error on a second access, and finalization scores the returned candidate and the seed on the same final tasks. The seal blocks scoring and editing but cannot stop reads. Plugin hooks, which live in core-owned scripts and not in editable skill text, deny edits to the sealed test and gold files and block a session from finishing until the adapter check and the gate are green.

Gold isolation has two layers. The adapter gives the target agent a gold-free copy of the task data and scores against the pristine copy. The proposer, a coding-agent subprocess with shell access, runs under a separate user id or in a mount namespace, and a negative test must show denial before launch. The library does not sandbox the proposer itself, so this is a deployment requirement, and a hash manifest over the evaluation surface can detect tampering but cannot stop a read. On the $\tau^2$ suites the evaluator replays the agent's tool calls on the pristine $\tau^2$ toolkit, so edits to tool code change how the agent acts but not how the reference state is computed. Section~\ref{sec:v2:setup} states what was isolated for each arm.

The shipped acceptance rule, the stock gate, is a paired test: accept a candidate when the mean per-task improvement over the current best exceeds $k=1$ standard errors of the paired differences, computed on the selection split only. Hill-climb first sets to zero the tasks whose rollouts never touch the edited files (the edit's footprint) and requires the candidate's own valid fraction to reach 0.6; agent-optimize judges with a $z$-test that has a standard-error floor. About half of all no-op edits measure as a small improvement by chance, so a gate at zero banks noise.

Rule P is the conservative alternative and is off by default. It has two stages. A candidate first passes a screen on the full, unmasked difference vector at $k=1$. It is then confirmed by re-evaluating candidate and incumbent on the same selection tasks with $n_{\mathrm{conf}}$ new trials per task and aligned seeds, and is promoted only if the confirmation improvement is large relative to its standard error ($z_{\mathrm{conf}}$) and enough tasks are covered ($\rho$). The confirmation, not the screen, becomes the next baseline, which removes the selection bias from the reference that later candidates must beat. We froze $n_{\mathrm{conf}}=3$, $z_{\mathrm{conf}}=1.282$, no minimum improvement and $\rho=0.95$ from a budget-exact simulation of one search per suite (Appendix~\ref{app:sim}); Algorithm~\ref{alg:P} gives the rule. Confirmation costs $2\,|D_{\mathrm{sel}}|\,n_{\mathrm{conf}}$ rollouts per screened candidate, charged to the budget.

Table~\ref{tab:sim} gives the operating characteristics, simulated at each suite's observed difficulty and gating-split size, for ten candidates per search, and an empirical check below tests them. In simulation, in an all-null search rule P promotes something falsely with probability \simPFalseA{} on Suite A, \simPFalseB{} on B and \simPFalseC{} on C. The one-stage comparator, an unmasked paired gate with no confirmation, does so with \simKoneFalseA{}, \simKoneFalseB{} and \simKoneFalseC{}; it is not the shipped stock gate, whose footprint mask the simulation omits. Rule P pays in power: a candidate with a true 8-point gain is promoted with probability \simPPowerEightA{}, \simPPowerEightB{} and \simPPowerEightC{}, against \simKonePowerEightA{}, \simKonePowerEightB{} and \simKonePowerEightC{}. Under a stress mix of 60\% null, 25\% positive and 15\% negative edits, the expected true gain of the final artifact is only +\simPGainA{}, +\simPGainB{} and +\simPGainC{} for rule P and +\simKoneGainA{}, +\simKoneGainB{} and +\simKoneGainC{} for the one-stage comparator, so neither exceeds half a point per search. An empirical check is stricter than the simulation. On 40 repeated evaluations of the unmodified Suite D seed (gpt-4.1-mini, 49 tasks), drawing null candidates from the pool, the one-stage unmasked rule promotes at least one of ten null candidates in \nullKone{} of searches (interval \nullKoneLo{} to \nullKoneHi{}), and rule P as run in \nullP{} (\nullPLo{} to \nullPHi{}). Rule P lowers false acceptance about \nullRatio{} times relative to the one-stage rule, but its absolute rate is about \nullPSimRatio{} times what the simulation predicted. The check covers unmasked rules and null edits on a different target from Suite B; it does not test the masked stock gate or true effects. Rule P does not control adaptivity over the task set: the selection tasks also feed the proposer, so a candidate can overfit them and pass both stages. Only the final set removes that bias, and only to the degree it was not reused.

Three further checks guard the measurement. A trial that errors because of the candidate scores zero, while an infrastructure error is retried and, if it persists, scores zero at final scoring (Section~\ref{sec:v2:stats}). The gate refuses to judge when too few tasks are valid on both sides of the pair. Null-edit controls and injected scorer and adapter faults test all of this. \needs{E04 null-edit false-accept rates and the E05 and E06 fault-injection outcomes (15 pass, 7 partial, 6 gap, 1 info at last count), with the gaps stated.}
